\section{Related Works}

\noindent\textbf{Geographic load balancing.}
Since the onset of cloud-based computing, geographic load balancing has been proposed as a way to balance compute across multiple regions towards lowering electricity prices and carbon intensity~\cite{qureshi2009cutting,liu2015greening}.
Within this GLB framing, work incorporates intermittency by coupling load shifting with renewable generation, storage, and backup-grid usage~\cite{liu2011geoloadbalance,goiri2011greenslot,berral2014buildinggreencloud}.
More broadly, “green cloud” control planes combine request routing, VM migration, admission control, and storage-aware scheduling to improve renewable utilization or reduce cost/carbon emissions~\cite{zhang2011greenware,radovanovic2023carbonaware, johnson2025signalawareworkloadshiftingalgorithms}.

\noindent\textbf{Grid-aware datacenter siting and placement.}
A smaller set of work directly couples datacenter siting decisions to power-system feasibility, moving beyond site-local signals to include transmission-network effects.
These models capture how placement choices interact with congestion, losses, and operational constraints (e.g., line overload and voltage violations)~\cite{goiri2011intelligentplacement,wang2015griddcwind}.
By embedding grid physics into the placement objective and constraints, this line motivates treating datacenters as grid-sited assets whose system-level impacts depend on \emph{where} capacity is added, not just on total demand.

\noindent\textbf{Datacenters for demand response.}
Complementary work studies datacenters as controllable demand capable of providing flexibility through policy-driven load modulation, workload shifting, and local generation/storage~\cite{irwin2011continuousdr,liu2013coincidentpeak}.
This literature primarily targets \emph{operational} coupling (minutes to hours), characterizing which forms of datacenter flexibility are practical, what signals they can respond to, and how these choices interact with utility cost structures and grid stress.

\noindent\textbf{Capacity Expansion Modeling Tools}
Grid-aware planning relies on optimization frameworks to explore new situations under uncertainty.
PyPSA and PyPSA-USA enable network-constrained, multi-period studies with renewables and storage~\cite{PyPSA,tehranchi5029120pypsa}, while capacity-expansion platforms such as SWITCH, Calliope, and oemof support multi-decade planning with high renewable penetration~\cite{nelson2012switch,johnston2019switch2,pfenninger2018calliope,hilpert2018oemof}.
Bilevel and expansion-planning formulations are well established in power systems~\cite{baringo2011wind}. Gradient-based methods can maker multi-scenario expansion planning scalable to larger problem sizes~\cite{degleris2024gradient,degleris2024gpu}.